\section{Introduction}

Autonomous-driving research already benefits from large image and multimodal corpora such as KITTI \cite{RefWorks:geiger2013vision}, Cityscapes \cite{RefWorks:cordts2016cityscapes}, Argoverse \cite{RefWorks:chang2019argoverse}, nuScenes \cite{RefWorks:caesar2020nuscenes}, Waymo Open \cite{RefWorks:sun2020scalability}, and BDD100K \cite{RefWorks:yu2020bdd100k}. These datasets are foundational for detection, segmentation, tracking, and motion forecasting, but they provide limited support for consistent, per-image scene semantics that can be queried across sources. This gap becomes especially visible when evaluating vision-language models (VLMs): object-centric labels and coarse metadata do not reveal where a model fails under specific contextual conditions such as degraded visibility, visible traffic control, or vulnerable road users at complex intersections.

Recent work on driving VLMs has focused on free-form question answering, instruction following, or direct scenario reasoning \cite{RefWorks:deruyttere2019talk2car,RefWorks:sima2024drivelm,RefWorks:qian2024nuscenesqa,RefWorks:ding2024holistic,RefWorks:rivera2025scenario}. Those efforts show that multimodal models can describe traffic scenes, but they do not provide a fixed, schema-aligned benchmark that makes scene-level failures easy to slice, compare, and reproduce across legacy autonomous-driving datasets. As a result, researchers who want to audit model behavior across weather, visibility, traffic control, or road geometry still need to assemble ad hoc metadata or free-text prompts on top of upstream datasets.

CARScenes addresses this problem as a data-centric dataset and benchmark rather than a control or policy benchmark. We package an annotation-only, schema-aligned release of 5,192 front-camera frames drawn from Argoverse1, Cityscapes, KITTI, and nuScenes. Each record describes scene type, time of day, weather, road conditions, lane structure, traffic signs and lights, surrounding vehicles, pedestrians, visibility, ego context, camera condition, and a coarse severity score. The release is canonicalized into a single JSONL format, validated against a machine-readable schema, and distributed with fixed split manifests, slice definitions, and executable benchmark tooling. This framing is deliberate: CARScenes is meant to support structured scene understanding and failure analysis for driving VLMs, not to stand in for a vision-language-action dataset.

The project repository, including code and release artifacts, is publicly available at \url{https://github.com/Croquembouche/CARScenes}.

The practical use case is failure slicing. CARScenes makes it possible to ask whether two models behave differently on scenes with adverse visibility and parked vehicles, on intersections containing vulnerable road users, or on traffic-light scenes in higher semantic-difficulty buckets. These are the units of analysis that matter for dataset auditing, scenario mining, and cross-dataset evaluation, yet they are not available in a consistent form from the source datasets themselves.

This paper makes three contributions. First, it turns four legacy driving datasets into a single validated annotation-only release with a stable schema, benchmark splits, and reusable tooling. Second, it documents a VLM-assisted but fully human-corrected annotation workflow, so the released benchmark records should be interpreted as human-reviewed ground truth rather than raw model output. Third, it positions CARScenes as infrastructure for scenario-focused evaluation and data-centric failure analysis, which is the practical contribution highlighted by the IV reviews and the most defensible fit for the NeurIPS Datasets and Benchmarks track.
